% =====================================================================
% 09_permutation_p123.tex
% 第 9 章 置换算符 P_{123} 的解析与代码实现
% Wave 3 / Agent A
%
% 配套代码（current）:
%   src/core/state\_space/make_permutation_matrix.{cpp,h}
%   src/utils/auxiliary.cpp        (pi1_tilde / pi2_tilde / Atilde / Gtilde)
%   src/io/disk_io_routines.cpp    (HDF5 cache)
%   src/core/faddeev_solver/solve_faddeev.cpp (P123 下游消费)
% 配套代码（origin）:
%   tictac-origin/Tic-tac/CPP/make_permutation_matrix.{cpp,h}
% =====================================================================

% --- 本章内部宏 (providecommand 防止与全局 physics_macros.tex 冲突) ---
\providecommand{\permmat}{\Pop}                              % 置换矩阵 P_{123}
\providecommand{\Pijkl}{P_{(12)(3)}}                          % 12 配对 → 3 配对
\providecommand{\jacP}{\boldsymbol{p}}                        % 12 相对动量
\providecommand{\jacQ}{\boldsymbol{q}}                        % 12-3 相对动量
\providecommand{\jacPp}{\boldsymbol{p}'}
\providecommand{\jacQp}{\boldsymbol{q}'}
\providecommand{\pibar}{\bar{\boldsymbol{\pi}}}
\providecommand{\pibarone}{\tilde{\pi}_{1}}
\providecommand{\pibartwo}{\tilde{\pi}_{2}}
\providecommand{\Atilde}{\tilde{A}}
\providecommand{\Gtilde}{\tilde{G}}
\providecommand{\binPp}[1]{\Delta^{p'}_{#1}}
\providecommand{\binQp}[1]{\Delta^{q'}_{#1}}
\providecommand{\NxGL}{N_{x}}
\providecommand{\Nphi}{N_{\phi}}
\providecommand{\Lmax}{L_{\max}}
\providecommand{\NJP}{N_{J^\pi}}
\providecommand{\nineJ}[9]{\left\{\begin{array}{ccc}#1&#2&#3\\#4&#5&#6\\#7&#8&#9\end{array}\right\}}
\providecommand{\sixJ}[6]{\left\{\begin{array}{ccc}#1&#2&#3\\#4&#5&#6\end{array}\right\}}
\providecommand{\threeJ}[6]{\begin{pmatrix}#1&#2&#3\\#4&#5&#6\end{pmatrix}}
\providecommand{\WPstate}[3]{\ket{X^{#1}_{#2#3}}}             % |X^α_{ij}>，i=q-bin、j=p-bin
\providecommand{\WPstateP}[3]{\ket{X^{#1}_{#2#3}}}            % bra/ket 同形

\chapter{置换算符 \texorpdfstring{$\Pop$}{P123} 的解析与代码实现}
\label{ch:09_permutation_p123}

% =====================================================================
\section{物理动机}
\label{sec:ch09-motivation}

% ---------------------------------------------------------------------
\subsection{从反对称化到 \texorpdfstring{$\Pop$}{P}}
\label{subsec:ch09-antisymmetrization}

三核子（3N）系统由两个相同的中子和一个质子（\(p+n+n\)）或两个质子和一个中子（\(d+p\)）组成，
即同位旋空间内三个核子全等。Pauli 反对称化要求所有物理态对任意两个核子的对换\emph{反对称}：
\begin{equation}
\label{eq:ch09-antisym}
\Aop\,\ket{\Psi}=\ket{\Psi},\qquad
\Aop \;=\; \tfrac{1}{6}\sum_{\sigma\in S_3}\sgn(\sigma)\,\hat{\sigma},
\end{equation}
其中 \(S_3\) 是三元置换群、\(\hat{\sigma}\) 把核子标号按 \(\sigma\) 重排。
但是 Faddeev 分解（\cref{ch:04_faddeev_ags}）只在\textbf{一种 Jacobi 树}上展开
（约定取 "(12)3" 树：\((12)\) 配对相互作用、\(3\) 旁观），因此每条 Faddeev 分量
\(\psifad{i}\) 都\textbf{打破}了三体置换对称性。要把这种"单一树"形式与 Pauli
反对称兼容，必须显式地把另两条树的内容\emph{借助置换算符搬运回来}。

记三个 Jacobi 树为 "(23)1"、"(31)2"、"(12)3"，对应的两两置换为 \(\Pij{1}{2}\)、
\(\Pij{2}{3}\)、\(\Pij{1}{3}\)。Glöckle 等人在 1983 年的标准论文（\cite{Gloeckle1983}）
中证明：
\begin{equation}
\label{eq:ch09-P-decomposition}
\Pop \;\equiv\; \Pij{1}{2}\Pij{2}{3} + \Pij{1}{3}\Pij{2}{3},
\end{equation}
便是把"(12)3"树映射到另两棵 Jacobi 树之和的\emph{唯一非平凡循环置换}。在反对称
\((12)\)-pair 子空间内（即 \(\Pij{1}{2}\ket{\Psi}=-\ket{\Psi}\)），\eqref{eq:ch09-P-decomposition}
还可以等价改写为
\begin{equation}
\label{eq:ch09-P-twoP123}
\Pop \;=\; 2\Pij{1}{2}\Pij{2}{3} \;\equiv\; 2\,\Pijkl,
\end{equation}
即下游求解器（\cref{ch:11_faddeev_solver}）实际遇到的"\(P\)"恰恰是单个三体循环
置换 \(P_{123}\equiv\Pij{1}{2}\Pij{2}{3}\) 的\textbf{两倍}。这一结论解释了
\texttt{solve\_faddeev.cpp} 内 \texttt{calculate\_PVC\_col} 函数中的
\verb|double P_element = 2*P123_val_array[idx_i]| 这一在初学者看来"莫名出现 2"
的代码段（详见 \cref{sec:ch09-interfaces}）。

\paragraph{命名约定。}
为避免读者混淆，本章后续\textbf{统一用 \(\Pop\) 指代}\eqref{eq:ch09-P-decomposition}
中的\textbf{完整置换算符}（factor 2 已经吸收进\(\Pop\) 的定义）；用
\(\Pijkl\) 或\(P_{123}\) 指代代码内存储的"半矩阵"。代码命名 \texttt{P123\_*}
对应后者；调用方负责乘 2。

% ---------------------------------------------------------------------
\subsection{为什么 \texorpdfstring{$\Pop$}{P} 决定 Faddeev 求解的精度}
\label{subsec:ch09-why-precision}

Faddeev/AGS 核（\cref{ch:04_faddeev_ags}）的真实操作形式是
\begin{equation}
\label{eq:ch09-faddeev-kernel}
T_{\alpha\beta}(z) \;=\; \delta_{\alpha\beta}\,\Vop_{\alpha}\,\Pop \;+\;
\Vop_{\alpha}\,\Pop\,\Gzero(z)\,T_{\alpha\beta}(z),
\end{equation}
其中 \(\Vop_\alpha\) 是配对 \(\alpha\) 上的两体核力（\cref{ch:08_potential_matrix}），
\(\Gzero(z)=(z-\Hzero)^{-1}\) 是自由 3N Green 函数（\cref{ch:10_resolvent}），
\(\Pop\) 是\eqref{eq:ch09-P-decomposition}所定义的置换。可以读出三件事：

\begin{enumerate}
\item Faddeev 核里的\textbf{每一次再散射}（rescattering）都\textbf{必须穿一次}
      \(\Pop\)。如果 \(\Pop\) 在 WP 基下的矩阵元有 \(\sim 10^{-3}\) 的相对误差，
      Padé 加速 12 阶后误差会放大到 \(\sim 10^{-2}\) 级，反映为 \(iT_{11}\)
      在 \(\theta_{\text{c.m.}}\sim 60^\circ\) 处的扭曲。
\item \(\Pop\) 是\textbf{与能量无关}的纯几何算符（仅由 \(\Hzero\)
      谱、Jacobi 坐标变换决定）。一旦 WP 基底固定（\(\NpWP, \NqWP\) 选定），
      所有 \(z\) 值的求解共享同一个 \(\Pop\)。这意味着\textbf{建一次、用许多次}
      ——这正是 \cref{sec:ch09-codeland} 看到的 HDF5 cache 设计动机。
\item \(\Pop\) 在 WP 基下\textbf{高度稀疏}。\eqref{eq:ch09-P-decomposition}
      把"(12)3"基矢映射到 \((23)1\) 与 \((31)2\) 的线性组合；几何上这只让
      "原 bin 与新 bin 重叠的格子"对中存在矩阵元。下面会量化：典型生产配置
      （\(\NpWP=\NqWP=20\)、\(\Nalpha=34\)、\(J^\pi=\tfrac{1}{2}^+\)）下，
      \(\Pop\) 的稠密维度为 \(13\,600^2\)，但非零元密度只有 \(\sim 0.5\%\)。
\end{enumerate}

\paragraph{结论：\(\Pop\) 是 Tic-tac 求解器最值得"花钱"的一次性投资。}
本章 \(\sim 90\%\) 的内容都在描述这一次性投资如何被分解为：（i）几何重叠
预筛、（ii）角动量 \(\Atilde\) 表预存、（iii）几何核 \(\Gtilde\) 表预存、
（iv）波包平均求积、（v）线程局部稀疏 COO 缓冲、（vi）HDF5 持久化与 mismatch
保护。这条流水线就是 \texttt{make\_permutation\_matrix.cpp} 这 1\,491 行代码
的物理含义。

% ---------------------------------------------------------------------
\subsection{上下游边界}
\label{subsec:ch09-up-down-stream}

为了在心智模型中先勾勒边界，先放出本模块的 I/O 表（具体的 TikZ 调用图见
\cref{sec:ch09-interfaces}）：

\begin{center}
\renewcommand{\arraystretch}{1.18}
\begin{tabular}{ll}
\toprule
\textbf{上游（输入）} & \\
\(\NpWP, \NqWP\) 与 bin 边界 \(\{p^{\text{WP}}_i\},\{q^{\text{WP}}_j\}\)
   & \cref{ch:06_wp_swp_states} 给出 \\
通道索引 \(\PWchan=(\ell, s, j; \lambda, I; J^\pi, T)\)
   & \cref{ch:07_pw_symm_states} 给出 \\
角动量上限 \(L_{2N}^{\max}, l_{1N}^{\max}, J_{2N}^{\max}\)
   & \texttt{run\_params} \\
求积阶 \(\NxGL\)（Gauss--Legendre）、\(\Nphi\)（极角网格）
   & 输入 \\
\midrule
\textbf{产出（输出）} & \\
COO 三元组 \((P_{123}^{\text{val}}, \text{row}, \text{col})\)
   & \texttt{P123\_sparse\_*\_array} \\
非零元数 \texttt{P123\_sparse\_dim} & \(\sim 10^{6}\) for \(\NpWP{=}20\) \\
HDF5 cache 文件 \texttt{P123\_sparse\_JP\_\dots.h5} & 供后续多能量复用 \\
\midrule
\textbf{下游（消费者）} & \\
\(\Pop\Vop\Cmat\)、\(\Cmat^T\Pop\Vop\Cmat\) 列向量构造
   & \cref{ch:11_faddeev_solver} \\
W1 三体力期望值 \(\bra{\psi_t}W^{(1)}\Pop\ket{\psi_t}\)
   & \cref{ch:17_3nf_numerical} \\
\bottomrule
\end{tabular}
\end{center}

% =====================================================================
\section{数学推导}
\label{sec:ch09-math}

% ---------------------------------------------------------------------
\subsection{矩阵元起步：\texorpdfstring{$\bra{\PWchanp\jacobip'\jacobiq'}\Pop\ket{\PWchan\jacobip\jacobiq}$}{<a'p'q'|P|apq>}}
\label{subsec:ch09-element-startpoint}

记 \cref{ch:03_jacobi_partial_wave} 中的部分波-Jacobi 基矢为
\(\ket{\PWchan; pq}\equiv \ket{\PWchan}\otimes \ket{p,q}\)，
其中 \(\PWchan=(L_{12},S_{12},J_{12},T_{12};l_3,s_3=\tfrac{1}{2};J^\pi,T)\)
是 \cref{eq:ch03-channel-tuple} 已建立的部分波元组。
\(p\) 是 \((12)\) 配对的相对动量、\(q\) 是 \(3\) 旁观核子相对于 \((12)\) 质心
的相对动量。在 \cref{ch:03_jacobi_partial_wave} 的 \eqref{eq:ch03-jacobi-12to23}
（一般性 Jacobi 旋转关系）下，置换 \(\Pij{2}{3}\) 把 "(12)3" 树的 \((p,q)\)
映射到 "(13)2" 树的 \((\bar{p},\bar{q})\)：
\begin{equation}
\label{eq:ch09-jacobi-rotation}
\bar{p} \;=\; \pibarone(p,q,x), \qquad
\bar{q} \;=\; \pibartwo(p,q,x), \qquad
x \;\equiv\; \hat{p}\!\cdot\!\hat{q},
\end{equation}
其中
\begin{equation}
\label{eq:ch09-pi1-pi2-explicit}
\pibarone(p,q,x) = \sqrt{\tfrac{1}{4}p^{2} + \tfrac{9}{16}q^{2} + \tfrac{3}{4}\,p\,q\,x}, \qquad
\pibartwo(p,q,x) = \sqrt{p^{2} + \tfrac{1}{4}q^{2} - p\,q\,x}.
\end{equation}
\eqref{eq:ch09-pi1-pi2-explicit}是 Glöckle 1983 \cite{Gloeckle1983} \S2 与
Sean Miller 博士论文 \cite{MillerThesis} \S6.4 都列出的"标准 Jacobi 旋转"，
也可以从如下三核子坐标推导（用核子 \(i\) 的位置/动量 \(\boldsymbol{r}_i,
\boldsymbol{k}_i\) 表达）：
\begin{align}
\boldsymbol{p}    &= \tfrac{1}{2}(\boldsymbol{k}_1 - \boldsymbol{k}_2), &
\boldsymbol{q}    &= \tfrac{2}{3}\bigl[\boldsymbol{k}_3 - \tfrac{1}{2}(\boldsymbol{k}_1+\boldsymbol{k}_2)\bigr],\\
\bar{\boldsymbol{p}} &= \tfrac{1}{2}(\boldsymbol{k}_1 - \boldsymbol{k}_3), &
\bar{\boldsymbol{q}} &= \tfrac{2}{3}\bigl[\boldsymbol{k}_2 - \tfrac{1}{2}(\boldsymbol{k}_1+\boldsymbol{k}_3)\bigr].
\end{align}
两组动量满足
\begin{equation}
\bar{\boldsymbol{p}} = -\tfrac{1}{2}\boldsymbol{p} - \tfrac{3}{4}\boldsymbol{q},\qquad
\bar{\boldsymbol{q}} = \boldsymbol{p} - \tfrac{1}{2}\boldsymbol{q}.
\end{equation}
取其模长平方再开方即得 \eqref{eq:ch09-pi1-pi2-explicit}。
代码端 \(\pibarone, \pibartwo\) 出现在 \texttt{src/utils/auxiliary.cpp:453--473}
的两个工具函数 \texttt{pi1\_tilde / pi2\_tilde}：

\lstinputlisting[style=cpp, caption={\(\pibarone, \pibartwo\) 的代码实现},
label={lst:ch09-pi-tilde}]
{code_excerpts/snippets/ch09_pi_tilde_helpers.cpp}

\begin{remark}[\(\tilde\pi'\) 的存在原因]
\label{rem:ch09-pi-prime}
\texttt{pi1\_prime\_tilde / pi2\_prime\_tilde} 与
\texttt{pi1\_tilde / pi2\_tilde} 仅相差 \(x\to-x\) 的整体符号；它们对应另一种
约定（\(\Pij{1}{3}\Pij{2}{3}\) vs \(\Pij{1}{2}\Pij{2}{3}\)）。见
\cref{sec:ch09-pitfalls} 关于符号约定的讨论。
\end{remark}

把 \eqref{eq:ch09-jacobi-rotation} 代回\(\bra{\PWchanp\jacobip'\jacobiq'}\Pop\ket{\PWchan\jacobip\jacobiq}\)，
利用动量本征态的归一化 \(\braket{\boldsymbol{p}'}{\boldsymbol{p}}=\delta(\boldsymbol{p}'-\boldsymbol{p})\)
做 \(\delta\) 函数的"双 Jacobi 套索"：
\begin{equation}
\label{eq:ch09-element-raw}
\bra{\PWchanp\jacobip'\jacobiq'}\Pop\ket{\PWchan\jacobip\jacobiq}
\;=\;
\int_{-1}^{+1}\!dx\;
\frac{\delta(p' - \pibarone(p,q,x))\,\delta(q' - \pibartwo(p,q,x))}{p'^{2}\,q'^{2}}\;
\Atilde_{\PWchanp\PWchan}(L_{\text{tot}},S_{\text{tot}})\,
\Gtilde_{\PWchan\PWchanp}(p,q,x).
\end{equation}
其中 \(\Atilde\) 是\textbf{角动量重耦合系数}（与 \(p,q\) 无关、只依赖通道 quantum
numbers）、\(\Gtilde\) 是\textbf{几何核}（依赖 \(p,q,x\) 的连续函数）。下面分别
推导。

% ---------------------------------------------------------------------
\subsection{角动量重耦合：\texorpdfstring{$\Atilde$}{Atilde} 的封闭形式}
\label{subsec:ch09-Atilde}

\(\Pop\) 把通道指标 \(\PWchan\to\PWchanp\) 的非平凡映射全部装进 \(\Atilde\)。
按 Glöckle \cite{Gloeckle1983} \S2.3，可以把每条通道的角动量耦合从
\((L_{12},l_3)\) 配对耦到 \((S_{12},s_3=\tfrac{1}{2})\) 配对、再耦到总
\((J,M_J)\)，并把"(12)3" 变 "(13)2" 的两次重耦合用 9-j 与 6-j 符号表达：
\begin{align}
\label{eq:ch09-Atilde-explicit}
\Atilde_{\alpha\alpha'}(L_{\text{tot}})
&=
\sum_{S_{\text{tot}}}\!
\sqrt{
\hat{J}_{12}\hat{j}_3\hat{S}_{12}\hat{T}_{12}
\hat{J}_{12}'\hat{j}_3'\hat{S}_{12}'\hat{T}_{12}'
}
\;(-1)^{S_{12}'+T_{12}'}\,(2 S_{\text{tot}}+1)\notag\\
&\quad\times
\nineJ
{2L_{12}}{2S_{12}}{2J_{12}}
{2l_{3}}{1}{2j_{3}}
{2L_{\text{tot}}}{2S_{\text{tot}}}{2J}\;
\nineJ
{2L_{12}'}{2S_{12}'}{2J_{12}'}
{2l_{3}'}{1}{2j_{3}'}
{2L_{\text{tot}}}{2S_{\text{tot}}}{2J}\notag\\
&\quad\times
\sixJ{1}{1}{2S_{12}}{1}{2 S_{\text{tot}}}{2S_{12}'}\;
\sixJ{1}{1}{2T_{12}}{1}{2 T}{2T_{12}'}.
\end{align}
其中我用了\(\hat{x}\equiv 2x+1\) 的简记法、且所有 \(j,J,T\) 均以
\(2j\) 整数形式表达（GSL 要求）。\(L_{\text{tot}}\) 是中间求和指标
（\(L_{12}\) 与 \(l_3\) 耦合的总轨道角动量），\(S_{\text{tot}}\)
则是 \(S_{12}\) 与 \(s_3=\tfrac{1}{2}\) 耦合的总自旋。两个 9-j 符号的物理
含义是\textbf{把对偶位置的 \((L_{12},S_{12},J_{12})\) 重耦到
\((L_{\text{tot}},S_{\text{tot}},J)\)}；两个 6-j 符号则负责\textbf{把
\(s_3\)、\(t_3\) 这种"半整数自旋"在 \((12)\) 配对里旋转}。

代码端 \(\Atilde\) 的实现是 \texttt{src/utils/auxiliary.cpp:863--884} 的
\texttt{Atilde()}：

\lstinputlisting[style=cpp, caption={\(\Atilde\) 与 \(\Gtilde\) 的代码实现},
label={lst:ch09-Atilde-Gtilde}]
{code_excerpts/snippets/ch09_atilde_gtilde_new.cpp}

\paragraph{Wigner 6j-/9j-符号的预存。}
\eqref{eq:ch09-Atilde-explicit}对每对通道 \((\alpha,\alpha')\)、对每个允许的
\(L_{\text{tot}}\)，需要查 4 个 Wigner 符号；外层 \(\PWchan\) 循环再叠加，
单纯 GSL 调用代价为 \(\mathcal{O}(\Nalpha^2 \Lmax)\)，但每次 GSL 调用本身约
\SI{1}{\micro\second}。Tic-tac 在
\texttt{make\_permutation\_matrix.cpp:466--491} 把所有
\((2j_1,\dots,2j_6)\le 2 l_{\max}\) 的 6-j 表预先用 \texttt{collapse(3)} OpenMP
循环填满，再用 inline \texttt{SixJSymbol(...)} 查表。同样的处理对 9-j 符号
（实际是\texttt{gsl\_sf\_coupling\_9j} 的 inline 调用）和 Clebsch--Gordan 系数
（\texttt{ClebschGordan\_data} 的 5-D 数组）一并执行。

\begin{remark}[Edmonds vs Brink 约定]
\eqref{eq:ch09-Atilde-explicit}的相位因子 \((-1)^{S_{12}'+T_{12}'}\) 用的是
Edmonds 约定（\cite{Edmonds1957}）。\cref{ch:07_pw_symm_states} 里
\texttt{make\_pw\_symm\_states.cpp} 也使用同一约定，因此与本章\(\Atilde\)兼容。
若读者从 Brink--Satchler \cite{BrinkSatchler1968} 或 Witała 等人的早期 Bochum
论文手抄角动量代数，注意 \((-1)^{T_{12}}\) vs \((-1)^{T_{12}'}\) 的差别；写错
就会让 W1 期望值整体相位翻转（\cref{ch:17_3nf_numerical} 中第 17.8 节有相应
踩坑记录）。
\end{remark}

% ---------------------------------------------------------------------
\subsection{几何核 \texorpdfstring{$\Gtilde$}{Gtilde} 的完整推导}
\label{subsec:ch09-Gtilde}

把 \eqref{eq:ch09-element-raw} 中的 \(\delta\) 函数对 \(p',q'\) 积掉，余下的
\(x\) 积分把通道之间的"动量+角动量"几何耦合压缩成一个标量函数
\(\Gtilde_{\alpha\alpha'}(p,q,x)\)。其封闭形式来自展开三个球谐 \(Y_{l m}\)：

\paragraph{第一步：球谐展开 \(\ket{p,q}\) 中的角部分。}
回忆 \cref{eq:ch03-angularY} 的部分波展开：
\begin{equation}
\ket{p,q;\PWchan}
=
\sum_{m_{L},m_{l}}
C^{J M_J}_{L_{12} m_L, l_3 m_l;\dots}
\,Y_{L_{12} m_L}(\hat{p})\,Y_{l_3 m_l}(\hat{q})\,
\ket{(\tfrac{1}{2}\tfrac{1}{2})S_{12};(S_{12}\tfrac{1}{2})S_{\text{tot}}}
\otimes\ket{p,q;\text{radial}},
\end{equation}
经 \(\Pij{2}{3}\) 的 Jacobi 旋转后 \((\hat{p},\hat{q})\to(\hat{\bar p},\hat{\bar q})\)，
两组方向相互"耦合到同一组动量"：使用球谐加法定理
\(Y_{Lm}(\hat{n})=Y_{Lm}(\theta(\hat{n}),\varphi(\hat{n}))\) 与
\(P^{m}_{L}(\cos\theta)\) 的解析表达式可以写出：
\begin{equation}
Y_{L_{12} m_L}(\hat{\bar p})\;=\;\sqrt{\tfrac{2L_{12}+1}{4\pi}}\,P^{m_L}_{L_{12}}(\cos\theta_1)\,
\delta_{m_L,0}\;\;\;\text{在 }\hat{p}=\hat{z}\text{ 取向},
\end{equation}
其中我们\textbf{选取}\(\hat{p}\) 沿 \(\hat{z}\) 方向（\(\hat{p}\) 是积分变量、可以
任选取向），\(\theta_1\) 是 \(\hat{\bar p}\) 与 \(\hat{p}\) 的夹角。从
\eqref{eq:ch09-pi1-pi2-explicit} 配合 \(\bar{\boldsymbol{p}}=-\tfrac{1}{2}\boldsymbol{p}-\tfrac{3}{4}\boldsymbol{q}\)
\(\Rightarrow\) \(\bar{p}\,\hat{\bar p}\!\cdot\!\hat{p}=-\tfrac{1}{2}p-\tfrac{3}{4}q\,(\hat{q}\!\cdot\!\hat{p})\)，
我们立即得到
\begin{equation}
\label{eq:ch09-costheta1}
\cos\theta_1 \;=\; -\frac{\tfrac{1}{2}p + \tfrac{3}{4}q\,x}{\pibarone(p,q,x)},
\qquad
\cos\theta_2 \;=\; \frac{p - \tfrac{1}{2}q\,x}{\pibartwo(p,q,x)}.
\end{equation}
代码端 \eqref{eq:ch09-costheta1} 出现在 \texttt{Gtilde\_new()} 函数
内（\texttt{src/utils/auxiliary.cpp:896--897}）以及 \texttt{calculate\_permutation\_elements\_for\_3N\_channel}
内对 Plm 表的预填段。两处都做了"夹紧到 \([-1,1]\)"的数值保护，因为
浮点累积误差会让 \(\cos\theta\) 偶尔达到 \(1+\eps\)。

\paragraph{第二步：把两个 \(\hat{\bar p}, \hat{\bar q}\) 的球谐合成一个"总球谐"。}
利用球谐耦合：
\begin{equation}
Y_{L_{12} m_L}(\hat{\bar p})\,Y_{l_3 m_l}(\hat{\bar q})
=
\sum_{L_{\text{tot}},M_{\text{tot}}}\!
C^{L_{\text{tot}} M_{\text{tot}}}_{L_{12} m_L, l_3 m_l}\,
\bigl[Y_{L_{12}}(\hat{\bar p})\otimes Y_{l_3}(\hat{\bar q})\bigr]_{L_{\text{tot}}M_{\text{tot}}}.
\end{equation}
配合 \(\eta\)-axis 取向使方位角 \(\varphi'=\pi\)（来自 \(\hat{\bar p}\) 与
\(\hat{\bar q}\) 的相对方位角约定：当 \(\hat{p}=\hat{z}\)、\(\hat{q}\) 在 \(xz\)-平面时，
\(\hat{\bar p}\) 与 \(\hat{\bar q}\) 共面但\textbf{反向}，导致 \(\varphi=\pi\)），
\(e^{i m\pi}=(-1)^m\) 给出 \(\Gtilde\) 表达式中的 \((-1)^{M_{\text{tot}}}\)
相位因子。

\paragraph{第三步：闭式的 \(\Gtilde\)。}
合并以上、把 \(\Atilde\) 的角动量耦合系数也吸进来，得：
\begin{empheq}[box=\fbox]{align}
\label{eq:ch09-Gtilde-final}
\Gtilde_{\alpha\alpha'}(p,q,x)
=
8\pi^2\!\!\sum_{L_{\text{tot}}=L_{\min}}^{L_{\text{max}}}\!\!
\Atilde_{\alpha\alpha'}(L_{\text{tot}})\!\!
\sum_{M=-l_3}^{l_3}\!\!
&C^{L_{\text{tot}}M}_{L_{12}\,0,\,l_3\,M}\,
\sqrt{\tfrac{2L_{12}+1}{4\pi}}\,(-1)^{M}\,P^{M}_{l_3}(x)\notag\\
\times\!\!\sum_{m=-L_{12}'}^{L_{12}'}\!\!
&C^{L_{\text{tot}}M}_{L_{12}'\,m,\,l_3'\,M-m}\,
P^{m}_{L_{12}'}(\cos\theta_1)\,P^{M-m}_{l_3'}(\cos\theta_2),
\end{empheq}
其中 \(L_{\min}=\max(\abs{L_{12}-l_3},\abs{L_{12}'-l_3'})\)、
\(L_{\max}=\min(L_{12}+l_3,L_{12}'+l_3',(2J+5)/2)\)；
\(C^{c\gamma}_{a\alpha,b\beta}\) 是 Clebsch--Gordan 系数（采用
\(2j\) 整数表示）。
\eqref{eq:ch09-Gtilde-final} 与 \eqref{eq:ch09-Atilde-explicit} 共同构成了
\eqref{eq:ch09-element-raw} 的右端被积函数。

\paragraph{与代码的对应。}
\eqref{eq:ch09-Gtilde-final} 与 \texttt{Gtilde\_new()}（\cref{lst:ch09-Atilde-Gtilde}）
是 1:1 对应：
\begin{itemize}
\item 第一个 \texttt{for (int Ltotal = ...)} 循环 = \(L_{\text{tot}}\) 求和；
\item \texttt{fac1 = 8.0 * M\_PI * M\_PI * A\_store[...]} = 前因子 \(8\pi^2\Atilde\)；
\item 第二个 \texttt{for (int Mtotal = ...)} 循环 = \(M\) 求和；
\item \texttt{ClebschGordan(2*L12, 2*l3, 2*Ltotal, 0, 2*Mtotal, 2*Mtotal)} = \(C^{L_{\text{tot}}M}_{L_{12}\,0,\,l_3\,M}\)；
\item \texttt{Plm(l3, Mtotal, x)} = \(P^{M}_{l_3}(x)\)；
\item 第三个 \texttt{for (int M12primesum = ...)} 循环 = \(m\) 求和（这里 \(m=M_{12}'\)）；
\item \texttt{Plm(L12prime, M12primesum, costheta1) * Plm(l3prime, Mtotal-M12primesum, costheta2)} =
      \eqref{eq:ch09-costheta1} 出来的两个 Legendre 项相乘。
\end{itemize}

\(\Gtilde\) 与 \(\Atilde\) 都是\textbf{与 WP bin 无关、只与连续 \(p,q,x\) 有关}
的纯几何对象——这正是后续把它们\emph{预存为表}的物理依据。

% ---------------------------------------------------------------------
\subsection{\texorpdfstring{$x\to(\phi,k)$}{x to phi-k} 的极坐标变换}
\label{subsec:ch09-polar-coordinates}

直接用 \eqref{eq:ch09-element-raw}做 \(\delta\)-函数积分有一个不愉快的副作用：
\(\delta(p'-\pibarone(p,q,x))\) 在 \(x\) 上的根 \(x_*(p',q',p,q)\) 是隐函数，
雅可比 \(\abs{\partial \pibarone/\partial x}\) 在 \(p\to 0\) 或 \(q\to 0\) 时
存在端点奇异。为了让积分光滑，Glöckle/Sean 都引入\textbf{极坐标变换}：把
\((p,q)\) 平面换成 \((\phi, k)\)：
\begin{equation}
\label{eq:ch09-polar}
p \;=\; k\sin\phi, \qquad q \;=\; k\cos\phi, \qquad
\phi\in(0,\tfrac{\pi}{2}),\quad k\in[0,\infty),
\end{equation}
在此变换下 \(\pibarone, \pibartwo\) 都\textbf{齐次}于 \(k\)：
\begin{equation}
\pibarone(k\sin\phi, k\cos\phi, x) = k\,\zeta_1(\phi,x),\qquad
\pibartwo(k\sin\phi, k\cos\phi, x) = k\,\zeta_2(\phi,x),
\end{equation}
其中
\begin{equation}
\zeta_1(\phi,x) = \pibarone(\sin\phi,\cos\phi,x),\qquad
\zeta_2(\phi,x) = \pibartwo(\sin\phi,\cos\phi,x).
\end{equation}
正是这个齐次性让 WP 平均（\cref{subsec:ch09-WP-avg}）的 \(k\) 积分变成
\textbf{解析的}（仅一个有理多项式 \((k_{\max}^2-k_{\min}^2)/2\zeta_1\zeta_2\)）；
留下的"难积分"只剩 \(\phi\) 与 \(x\) 两维，刚好可以用 Gauss--Legendre 求积。

代码端 \eqref{eq:ch09-polar} 出现在 \texttt{calculate\_WP\_overlap()}
（\texttt{src/core/state\_space/make\_permutation\_matrix.cpp:205--365}）的
"\texttt{double phi\_lower = atan(pp\_l/qp\_u); double phi\_upper = atan(pp\_u/qp\_l);}"
两行——这两行界定了\textbf{对每对 \((q'_{\text{bin}}, p'_{\text{bin}})\)}
的 \(\phi\) 积分上下限：
\begin{equation}
\phi\in[\arctan(p'_{\text{lo}}/q'_{\text{hi}}),\;\arctan(p'_{\text{hi}}/q'_{\text{lo}})].
\end{equation}
\(\Nphi\) 个 Gauss 节点会自动落在这个区间内（见 \cref{lst:ch09-WP-overlap}）。

% =====================================================================
\section{离散化方案}
\label{sec:ch09-discretization}

% ---------------------------------------------------------------------
\subsection{波包（WP）平均}
\label{subsec:ch09-WP-avg}

把 \eqref{eq:ch09-element-raw} 投影到 WP 基：每个 WP 基矢
\(\WPstate{\PWchan}{j}{i}\) 由 \cref{ch:06_wp_swp_states} 的两体 wave-packet
\(\ket{x_i^p}\) 与三体 wave-packet \(\ket{x_j^q}\) 张成，归一化为
\begin{equation}
\WPstate{\PWchan}{j}{i} \;\equiv\;
\frac{1}{\sqrt{\Niso^p_i\Niso^q_j}}
\int_{p_i^{\text{lo}}}^{p_i^{\text{hi}}}\!dp\;p\;
\int_{q_j^{\text{lo}}}^{q_j^{\text{hi}}}\!dq\;q\;
\ket{\PWchan; pq},
\end{equation}
其中 \(\Niso^p_i\equiv \int_{p_i^{\text{lo}}}^{p_i^{\text{hi}}}p^2 dp = (p_i^{\text{hi},3}-p_i^{\text{lo},3})/3\)，
\(\Niso^q_j\) 同理（\cref{ch:06_wp_swp_states} \eqref{eq:ch06-WP-norm}）。

WP 矩阵元定义为
\begin{equation}
\label{eq:ch09-WP-element}
P_{\PWchanp j' i', \PWchan j i}
\;\equiv\;
\bra{X^{\PWchanp}_{j' i'}}\Pop\ket{X^{\PWchan}_{j i}}.
\end{equation}
把 \eqref{eq:ch09-element-raw} 代入并先做 \(p',q'\) 的两个 WP 积分：
\(\delta(p'-\pibarone)\) 把 \(p'\) 钉死在 \(\pibarone\)；类似地 \(q'\) 钉死在
\(\pibartwo\)。\(p',q'\) 的 WP 积分把这两个 \(\delta\) 转化为 Heaviside 阶跃：
\(\Theta(\pibarone\in[\binPp{i'}])\,\Theta(\pibartwo\in[\binQp{j'}])\)。

接着对 \(p,q\) 的 WP 积分换到 \((\phi,k)\)（\eqref{eq:ch09-polar}），
\(dp\,dq=k\,dk\,d\phi\)。把所有 Heaviside 化为 \(k\) 的上下限：
\begin{align}
k_{\min}(\phi,x)
&= \max\!\left(
\tfrac{p_i^{\text{lo}}}{\sin\phi},\;\tfrac{q_j^{\text{lo}}}{\cos\phi},\;
\tfrac{p_{i'}^{\text{lo}}}{\zeta_1(\phi,x)},\;\tfrac{q_{j'}^{\text{lo}}}{\zeta_2(\phi,x)}
\right),\notag\\
k_{\max}(\phi,x)
&= \min\!\left(
\tfrac{p_i^{\text{hi}}}{\sin\phi},\;\tfrac{q_j^{\text{hi}}}{\cos\phi},\;
\tfrac{p_{i'}^{\text{hi}}}{\zeta_1(\phi,x)},\;\tfrac{q_{j'}^{\text{hi}}}{\zeta_2(\phi,x)}
\right).
\label{eq:ch09-kminmax}
\end{align}
\eqref{eq:ch09-kminmax}恰是\(\zeta_1\zeta_2\)、\(\sin\phi\cos\phi\)
等正定量的"四个比值取 max/min"——代码 \cref{lst:ch09-element-quad}
里的 \verb|kmin = max(kmin_array, kmin_array+4)|
\verb|kmax = min(kmax_array, kmax_array+4)| 即此。

\(k\) 积分变成
\begin{equation}
\int_{k_{\min}}^{k_{\max}}\!dk\;k\;\Bigl[\frac{1}{p'^{2}q'^{2}}\;\delta\text{两次}\Bigr]
\;\xrightarrow[\text{解析}]{}\;
\frac{k_{\max}^{2}-k_{\min}^{2}}{2\,\zeta_1(\phi,x)\zeta_2(\phi,x)}.
\end{equation}

合并所有，得到\textbf{WP 矩阵元的最终量化形式}：
\begin{empheq}[box=\fbox]{equation}
\label{eq:ch09-final-quadrature}
P_{\PWchanp j' i', \PWchan j i}
=
\frac{1}{\sqrt{\Niso^p_i\Niso^q_j\Niso^p_{i'}\Niso^q_{j'}}}\!
\int_{x=-1}^{+1}\!\!dx\!\int_\phi\! d\phi\;
\sin\phi\cos\phi\;
\Gtilde_{\alpha\alpha'}(\sin\phi,\cos\phi,x)\;
\frac{k_{\max}^{2}-k_{\min}^{2}}{2\,\zeta_1\zeta_2}.
\end{empheq}
\eqref{eq:ch09-final-quadrature} 是 1\,491 行代码所要数值化的"那一个公式"。

% ---------------------------------------------------------------------
\subsection{求积阶选择}
\label{subsec:ch09-quadrature-orders}

\eqref{eq:ch09-final-quadrature}中两个 1-D 求积：

\paragraph{\(x\) 方向：Gauss--Legendre 阶 \(\NxGL\)。}
\(x=\hat p\!\cdot\!\hat q\in[-1,+1]\)。被积函数关于 \(x\) 是
\((L_{12}+l_3+L_{\text{tot}}+5)\) 阶左右的多项式（来自
\(P^M_l(x)\) 与 \(\cos\theta_{1,2}\) 中 \(x\) 的依赖）。生产配置 \(\NxGL=8\)
对 \(L_{12}\le 4, l_3\le 4\) 完全积准；增至 \(\NxGL=16\) 验证一次后差别
\(<10^{-9}\)。Tic-tac 用 \(\NxGL=8\)（\texttt{run\_params.Nx}）作为默认。

\paragraph{\(\phi\) 方向：Gauss--Legendre 阶 \(\Nphi\)，区间为
\(\phi\in[\arctan(p_{i'}^{\text{lo}}/q_{j'}^{\text{hi}}),\arctan(p_{i'}^{\text{hi}}/q_{j'}^{\text{lo}})]\)。}
\textbf{此区间是\((q'_{\text{bin}}, p'_{\text{bin}})\)依赖的}！这是 Tic-tac 与
naive 数值求积的一个关键差别：被 \texttt{calculate\_WP\_overlap()} 收集到一个
\(\NqWP\!\times\!\NpWP\!\times\!\Nphi\) 的 \(\phi\) 节点表，索引方式见
\cref{lst:ch09-indexing-helpers} 的 \texttt{phi\_packet\_index} 内联函数。
生产配置 \(\Nphi=8\)；同样验证 \(\Nphi=16\) 与 \(\Nphi=8\) 差别 \(<10^{-7}\)。

\paragraph{\(k\) 方向：解析。}
如 \eqref{eq:ch09-final-quadrature}，\(k\) 积分被解析消掉，没有求积阶选择。

% ---------------------------------------------------------------------
\subsection{结构性稀疏：几何重叠预筛}
\label{subsec:ch09-WP-overlap-mask}

并不是每对 WP bin \(((q'_{\text{bin}}, p'_{\text{bin}}), (q_{\text{bin}}, p_{\text{bin}}))\)
都有非零矩阵元——只有当 \(k_{\min}<k_{\max}\) 在某个 \((\phi,x)\) 上成立时，
\eqref{eq:ch09-final-quadrature} 的被积函数才不全为零。Tic-tac 在做任何昂贵
求积\textbf{之前}先扫一遍这条不等式，把所有"几何不重叠"的 bin 配对从计算
循环中拿掉——把全 \(\NpWP^2\NqWP^2\) 网格压到大约 \(0.5\%\) 的非零密度。

代码端这一步是 \texttt{calculate\_WP\_overlap()}
（\texttt{src/core/state\_space/make\_permutation\_matrix.cpp:205--365}），
其布尔输出 \texttt{pq\_WP\_overlap\_array[$\NqWP^2\NpWP^2$]} 把允许的 bin 配
对编码为 1。这相当于一次 4-D
布尔下三角扫描，复杂度 \(\mathcal{O}(\NqWP^2\NpWP^2\Nphi\NxGL)\)，但常数极小。
注意这次扫描\textbf{不依赖通道指标}，因此整个 \(J^\pi\) 块共享同一张
overlap 表。

% ---------------------------------------------------------------------
\subsection{\texorpdfstring{$J^\pi$}{Jpi} 块对角化简}
\label{subsec:ch09-Jpi-block}

\(\Pop\) 与三体总角动量 \(\boldsymbol{J}\)、宇称 \(\Pi\)、同位旋第三分量
\(T_z\) 都对易，因此在 \((J^\pi,T)\)-标记的部分波分组下\textbf{严格块对角}。
Tic-tac 的 \texttt{run\_organizer} 主流程对每一个 \((J^\pi,T)\) 通道分别建一次
\(\Pop\)、求一次 Faddeev、再汇总观测量。第 9 章的所有"P\_123 矩阵"都默认
是单 \(J^\pi\) 块。下游\(\NJP\) 个块的总规模等于
\(\sum_{J^\pi}\Nalpha(J^\pi)\NpWP\NqWP\)（典型 \(\sim 10^4\sim 10^5\)）。

% =====================================================================
\section{算法骨架}
\label{sec:ch09-algorithm}

整个 \(\Pop\) 构造分四个抽象层。最外层入口 \texttt{fill\_P123\_arrays()}
负责"建一次"还是"读一次"；其余三层是计算细节。

\begin{algorithm}[H]
\DontPrintSemicolon
\caption{\textsc{Build-P123} 主流程}
\label{alg:ch09-build-P123}
\KwIn{\(\NpWP, \NqWP\)、bin 边界数组、通道索引 \(\PWchan\)、求积阶
       \(\NxGL,\Nphi\)、并行线程数 \(\Theta\)、HDF5 cache 路径}
\KwOut{COO 三元组 \((P_{123}^{\text{val}}, \text{row}, \text{col})\)
       与非零元数；HDF5 持久化文件}
\BlankLine
\If{\textnormal{\texttt{run\_params.calculate\_and\_store\_P123}}}{
  \tcp{Stage 1: 几何重叠预筛 (\texttt{calculate\_WP\_overlap})}
  分配 \(\NpWP^2\NqWP^2\) 布尔数组 \texttt{pq\_WP\_overlap}\;
  \ForEach{\(\phi\)-mesh \((q'_{\text{bin}}, p'_{\text{bin}})\)}{
    Gauss--Legendre 节点填到 \(\phi\)-array\;
    \ForEach{\((q_{\text{bin}}, p_{\text{bin}}, \phi, x)\) 4-D 网格}{
      检查 \(k_{\min}<k_{\max}\)，标记 overlap\;
    }
  }
  \tcp{Stage 2: 角动量与几何前预存
        (\texttt{calculate\_permutation\_elements\_for\_3N\_channel}, prelude)}
  预存 6-j 表 \(\mathcal{O}((2\Lmax+1)^6)\)\;
  预存 \(\Atilde_{\alpha\alpha'}(L_{\text{tot}})\) 表 \(\mathcal{O}(\Nalpha^2\Lmax)\)\;
  预存 Clebsch--Gordan 5-D 表\;
  预存 Plm 表（\(L_{12}, l_3\) 各一张，\(\Nphi\NxGL\NqWP\NpWP\) 项）\;
  \tcp{Stage 3: 主循环（\(\Pop\)填值），OpenMP \(\Theta\) 线程}
  \texttt{omp\_set\_num\_threads}\((\Theta)\)\;
  \ForEach{$(q'_{\text{bin}}, p'_{\text{bin}})$ 行 WP 块（外层 OpenMP）}{
    \ForEach{$(\alpha', \alpha)$ 通道对}{
      构造 \(\Gtilde_{\alpha\alpha'}\) 子表（\(\Nphi\!\times\!\NxGL\)）\;
      \ForEach{$(q_{\text{bin}}, p_{\text{bin}})$ 列 WP 块}{
        \If{\textnormal{\texttt{pq\_WP\_overlap}} 为真}{
          调用 \(\Gtilde\)+\(\zeta_{1,2}\) 计算 \eqref{eq:ch09-final-quadrature}\;
          \If{$|P_{123}|>0$}{
            写入线程本地 COO 缓冲\;
            \If{\textnormal{缓冲达 1 GB}}{
              flush 到 \texttt{P123\_subsparse\_*\_TFC\_*.h5}\;
              \texttt{TFC++}\;
            }
          }
        }
      }
    }
  }
  \tcp{Stage 4: 合并 (\texttt{read\_and\_merge\_thread\_files})}
  从 \(\Theta\) 个线程的所有 TFC 文件读回 COO 片段\;
  按 row-major 索引排序到稳定 COO\;
  \tcp{Stage 5: 持久化 (\texttt{store\_sparse\_permutation\_matrix\_for\_3N\_channel\_h5})}
  写入 \texttt{P123\_sparse\_JP\_*.h5}（含 bin 边界 + PW 元组）\;
}
\Else{
  \tcp{Recovery 路径：从 HDF5 cache 读回}
  \texttt{read\_sparse\_permutation\_matrix\_for\_3N\_channel\_h5}\;
  与当前运行的 \(\NpWP, \NqWP\)、bin 边界、PW 元组逐项对比，不匹配则
  \texttt{raise\_error}\;
}
\end{algorithm}

\paragraph{复杂度。}
\begin{itemize}
\item \textbf{时间}：主循环为
      \(\mathcal{O}\bigl(\Nalpha^2\NpWP^2\NqWP^2\Nphi\NxGL\bigr)\)。
      生产配置 \(\Nalpha=34\)、\(\NpWP=\NqWP=20\)、\(\Nphi=\NxGL=8\) 时
      原始操作数约 \(7.4\times 10^{9}\)；OpenMP \(\Theta=16\) 加速后
      墙钟约 \SI{15}{\minute} 到 \SI{30}{\minute}。
\item \textbf{峰值内存}：Plm 表
      \(\sim (\Nphi\NxGL\NqWP\NpWP)\!\times\!(L_{12}\!+\!l_3\!+\!1)^2\!\times\!8\,\textrm{B}\)
      \(\approx \SI{200}{\mega\byte}\)；Clebsch--Gordan 5-D 表 \(\sim
      \SI{100}{\mega\byte}\)；6-j 6-D 表 \(\sim \SI{10}{\mega\byte}\)；
      线程本地 COO 缓冲 \(\Theta\!\times\!\SI{1}{\giga\byte}\)。
\item \textbf{磁盘}：每个 \((J^\pi,T)\) 块的 HDF5 cache
      \(\sim \SI{50}{\mega\byte}\)；典型 18 个块累计 \(\sim
      \SI{1}{\giga\byte}\)。
\end{itemize}

% =====================================================================
\section{代码落地（origin / current 双栏对照）}
\label{sec:ch09-codeland}

本节给出 4 段代码——索引 helper、几何重叠、主循环、HDF5 cache——并对每段
逐行注释关键差异。所有片段都通过 \texttt{code\_excerpts/extract.py} 真实
抽取，保留原始行号。

% ---------------------------------------------------------------------
\subsection{代码片段 1：索引 helper（current 仓库新增）}
\label{subsec:ch09-code-indexing}

origin 仓库把所有 \texttt{packet\_state\_index = alpha*Nq*Np + q*Np + p}
直接展开在主循环里；current 仓库（commit 主分支）把它抽成三个 inline 函数
保护在匿名命名空间。这是\textbf{重构 diff 中最频繁出现的微变化}——origin
里"\(\alpha\NqWP\NpWP+q\NpWP+p\)" 散落 10 次，current 全部用
\texttt{packet\_state\_index} 替换。

\lstinputlisting[style=cpp,
  caption={current: 三个索引 helper inline 函数（packet/overlap/phi）},
  label={lst:ch09-indexing-helpers}]
{code_excerpts/snippets/ch09_indexing_helpers_current.cpp}

\begin{itemize}
\item \texttt{packet\_state\_index}: 三体 packet basis 的稠密索引
      \((\alpha,q,p)\to\alpha\NqWP\NpWP+q\NpWP+p\)。\(\Pop\) 的 row/col 都用
      此函数算得；这一致性保证了后续 \texttt{solve\_faddeev} 的 PVC 列向
      量索引（\cref{lst:ch09-PVC-consumer}）能直接对上号。
\item \texttt{wp\_overlap\_index}: \texttt{pq\_WP\_overlap} 数组的 4-D
      索引。注意它的 stride 是 \emph{列优先}（\(q'\) 最外、\(p\) 最内），
      与 \texttt{packet\_state\_index} 的 row 索引顺序一致。
\item \texttt{phi\_packet\_index}: \(\phi\)-mesh 的 3-D 索引。每对
      \((q'_{\text{bin}}, p'_{\text{bin}})\) 各 \(\Nphi\) 个 Gauss--Legendre
      节点。
\end{itemize}

\paragraph{重构动机：可读性 + 防错。}
当 \texttt{Nq\_WP*Np\_WP*Nphi} 这种乘积写到 5 处以上时，任何 \texttt{Nq\_WP}
误写为 \texttt{Np\_WP} 或漏掉一个 \texttt{Nphi} 都会导致越界。Sean 在
origin 仓库已经踩过两次这种坑（git log 里 commit 都带"index fix"）。重构后
所有越界都可以被 ASAN 立即检出，而且 IDE 的 "Find References" 能直接列出所
有调用点。

% ---------------------------------------------------------------------
\subsection{代码片段 2：几何重叠预筛 \texttt{calculate\_WP\_overlap}}
\label{subsec:ch09-code-WP-overlap}

下面把\textbf{current 与 origin 的同段代码并排呈现}（main loop 部分）。
相同：算法整体结构（外层 \texttt{omp parallel for}、4-D
循环、\(\phi\)-mesh 内嵌生成）；不同：current 用 \texttt{phi\_packet\_index}
helper 抽取索引、用 \texttt{wp\_overlap\_index} helper 抽取布尔数组下标，
而 origin 全部展开。

\noindent\textbf{Current 仓库}（\texttt{src/core/state\_space/make\_permutation\_matrix.cpp:205--365}）

\lstinputlisting[style=cpp,
  caption={current: \texttt{calculate\_WP\_overlap} 完整实现},
  label={lst:ch09-WP-overlap}]
{code_excerpts/snippets/ch09_calculate_WP_overlap_current.cpp}

\noindent\textbf{Origin 仓库}（\texttt{tictac-origin/Tic-tac/CPP/make\_permutation\_matrix.cpp:159--322}）

\lstinputlisting[style=cpp,
  caption={origin: \texttt{calculate\_WP\_overlap} 完整实现},
  label={lst:ch09-WP-overlap-origin}]
{code_excerpts/snippets/ch09_calculate_WP_overlap_origin.cpp}

\paragraph{逐行差异点。}
\begin{enumerate}
\item \textbf{\(\phi\)-mesh 索引}：origin 用裸表达式
      \verb|&phi_array[(qp_idx_WP*Np_WP + pp_idx_WP)*Nphi]|，
      current 用 \verb|&phi_array[phi_base_idx]| 其中
      \verb|phi_base_idx = phi_packet_index(qp_idx_WP, pp_idx_WP, 0, Np_WP, Nphi)|。
      数学语义完全等价。
\item \textbf{布尔数组索引}：origin 用 4 个 \texttt{step\_length\_*} 局部
      变量 + 内联展开；current 用 \verb|wp_overlap_index(...)| 调用。
\item \textbf{算法分支}：两段代码都保留了 \texttt{if (true) ... else \{ ...\}}
      的 dead-code 分支——后者是 Sean 论文里"naive 边界判断"的实现，
      没正确处理 \(\phi\) 上下限的极限情况，被 \texttt{if (true)} 永久禁用，
      但代码仍保留作为参考实现。current 与 origin 同此设计。
\item \textbf{OpenMP 并行}：两段都对最外层 \texttt{qp\_idx\_WP} 做
      \texttt{omp for}。原因：\(\phi\)-mesh 的初始化是 \texttt{calc\_gauss\_points}
      调用，对每对 \((q'_{\text{bin}}, p'_{\text{bin}})\) 写
      \(\Nphi\) 个连续位置；以 \texttt{qp} 切分线程可避免假共享。
\end{enumerate}

% ---------------------------------------------------------------------
\subsection{代码片段 3：主循环（构造矩阵元）}
\label{subsec:ch09-code-main-loop}

主循环的"骨架代码"约 \(150\) 行。现在并排展示 current 与 origin 的关键
\(\Pop\)填值循环（即从 row WP 外层、到 \texttt{calculate\_P123\_element\_in\_WP\_basis\_mod}
内层调用、再到 TFC flush 的全过程）。

\noindent\textbf{Current 仓库}（\texttt{src/core/state\_space/make\_permutation\_matrix.cpp:742--956}）

\lstinputlisting[style=cpp,
  caption={current: 主循环 + Gtilde 子表 + TFC flush},
  label={lst:ch09-main-loop}]
{code_excerpts/snippets/ch09_pe3n_main_loop_current.cpp}

\noindent\textbf{Origin 仓库}（\texttt{tictac-origin/Tic-tac/CPP/make\_permutation\_matrix.cpp:720--886}）

\lstinputlisting[style=cpp,
  caption={origin: 主循环 + Gtilde 子表 + TFC flush},
  label={lst:ch09-main-loop-origin}]
{code_excerpts/snippets/ch09_pe3n_main_loop_origin.cpp}

\paragraph{逐行差异点。}
\begin{enumerate}
\item \textbf{Gtilde 缓冲}：origin 用 VLA \verb|double Gtilde_subarray[Gtilde_subarray_size]|
      （C99 扩展，gcc 默认接受、clang 警告）；current 用 \verb|std::vector<double> Gtilde_subarray(...)|
      并 \verb|.data()| 取裸指针——更安全（不会 stack overflow），代价是一
      次堆分配/线程，对热路径无影响。
\item \textbf{Pakcet 索引}：origin 内联
      \verb|alphap_idx*Nq_WP*Np_WP + qp_idx_WP*Np_WP + pp_idx_WP|；
      current 一律 \verb|packet_state_index(alphap_idx, qp_idx_WP, pp_idx_WP, ...)|。
\item \textbf{P123 row/col 顺序}：两边\textbf{相同}——row 用
      \((\alpha', q', p')\)、col 用 \((\alpha, q, p)\)。这一致性是后续
      sort + CSC 转换的前提（\cref{subsec:ch09-code-merge}）。
\item \textbf{Pre-storage 头部}：origin 没有 inline 注释；current 加了大量
      \texttt{[EN]/[CN]} 双语注释解释每段的物理含义（详见
      \cref{lst:ch09-prelude} 的开头）。这是 current 重构最显著的"文档增厚"。
\end{enumerate}

\begin{lstlisting}[style=cpp, caption={current 主例程头部的双语注释（节选自 prelude）},
                   label={lst:ch09-prelude}]
// [EN] This routine is the expensive one-time construction of the
// finite-dimensional permutation matrix P. In the Miller paper and in
// chapter 6 of the notes, this is the step that replaces the continuous
// Jacobi-coordinate transformation by a reusable sparse matrix acting
// on packet states.
// [CN] 这个例程负责一次性构造有限维置换矩阵 P；按 Miller 论文和第 6 章
// 讲稿的说法，这一步就是把连续的 Jacobi 坐标变换替换成一个可复用的
// 稀疏波包矩阵。
\end{lstlisting}

\paragraph{核心被积函数 \texttt{calculate\_P123\_element\_in\_WP\_basis\_mod}。}
\eqref{eq:ch09-final-quadrature}的最终求积代码：

\lstinputlisting[style=cpp,
  caption={\(\Pop\) 的 WP 基矩阵元求积器（\texttt{src/utils/auxiliary.cpp:660--746}）},
  label={lst:ch09-element-quad}]
{code_excerpts/snippets/ch09_p123_element_quadrature.cpp}

\paragraph{逐行映射。}
\begin{itemize}
\item Line 21--32（current 抽取后行号）：bin 边界变量名 \(p_l, p_u, q_l, q_u, p'_l, p'_u, q'_l, q'_u\)
      ——\eqref{eq:ch09-kminmax} 中的 4 对端点。
\item Line 50--80：\(x\) 与 \(\phi\) 的双重 Gauss--Legendre 求积。
      \texttt{zeta\_1 = pi1\_tilde(sin\_phi, cos\_phi, x)} 即 \(\zeta_1(\phi,x)\)。
      \texttt{kmin\_array} / \texttt{kmax\_array} 直接写出 \eqref{eq:ch09-kminmax}。
\item Line 83--85：\eqref{eq:ch09-final-quadrature}的被积积子\
      \verb|wx * wphi * cos_phi*sin_phi * Gtilde * (kmax^2-kmin^2) / (2*zeta_1*zeta_2)|。
\item Line 90--94：除以\(\sqrt{\Niso^p_i\Niso^q_j\Niso^p_{i'}\Niso^q_{j'}}\)
      四个 WP 归一化常数（\texttt{p\_normalization}, \texttt{q\_normalization}
      在 \cref{ch:06_wp_swp_states} 已定义）。
\end{itemize}

% ---------------------------------------------------------------------
\subsection{代码片段 4：HDF5 持久化与稳健性检查}
\label{subsec:ch09-code-h5cache}

\(\Pop\) 与能量无关，因此一旦在某个 \(\NpWP, \NqWP\) 配置下建好，可以缓存
为 HDF5 文件、被多次复用。但是\textbf{cache 必须配套验证}——如果用户改了
\(\NpWP\) 但忘了清 cache，求解器会读到错误尺寸的矩阵导致段错误或更糟的
\textbf{静默错误}。Tic-tac 在读 cache 时立刻和当前运行的 \(\NpWP, \NqWP\)、
bin 边界数组、PW 元组逐项 \texttt{raise\_error}。

\lstinputlisting[style=cpp,
  caption={current: HDF5 写入 \(\Pop\) cache 主体（\texttt{src/io/disk\_io\_routines.cpp:1078--1165}）},
  label={lst:ch09-h5-store}]
{code_excerpts/snippets/ch09_h5_store_routine.cpp}

\lstinputlisting[style=cpp,
  caption={current: HDF5 读 cache 时的稳健性验证（\texttt{src/io/disk\_io\_routines.cpp:1209--1284}）},
  label={lst:ch09-h5-read}]
{code_excerpts/snippets/ch09_h5_read_validation.cpp}

\paragraph{文件命名规则}（\texttt{src/core/state\_space/make\_permutation\_matrix.cpp:185--199}）：

\lstinputlisting[style=cpp,
  caption={线程文件命名规则},
  label={lst:ch09-filename-rule}]
{code_excerpts/snippets/ch09_subarray_filename.cpp}

\paragraph{Cache key 包含的元数据。}
\begin{itemize}
\item \(2J^{(3N)}, P^{(3N)}\)：决定 \(J^\pi\) 块；
\item \(\NpWP, \NqWP\)：决定矩阵尺寸；
\item \(J_{2N}^{\max}\)：决定 \(\Nalpha\)；
\item Thread index + TFC：合并阶段文件区分。
\end{itemize}

合并到最终 \texttt{P123\_sparse\_JP\_$\dots$.h5} 时，Cache 内还写入：
\(\Nalpha\)、\(\NpWP, \NqWP\)、bin 边界数组、PW 元组数组、矩阵元 COO 三元组。
读回时\textbf{逐项校验}。

% ---------------------------------------------------------------------
\subsection{代码片段 5：合并与 COO 排序}
\label{subsec:ch09-code-merge}

OpenMP 各线程独立 flush 出来的 TFC 文件\textbf{无序}（每个线程
按 row WP 外循环、但线程间任意交错）。下游求解器要求 row-major COO
（即 row 索引非递减、行内 col 索引非递减），因此合并后必须再排序一次。

\lstinputlisting[style=cpp,
  caption={current: TFC 合并 + COO 全局排序（\texttt{make\_permutation\_matrix.cpp:1015--1220}）},
  label={lst:ch09-merge}]
{code_excerpts/snippets/ch09_merge_thread_files_current.cpp}

\paragraph{排序代价。}
\(\NJP\!\times\!\sim\!10^6\) 个 COO 三元组、单线程 \texttt{sort\_indexes}
（\texttt{utils/templates.h} 模板）大约 \SI{30}{\second}。如果换到
\texttt{std::sort + lambda} 仍是 \(\mathcal{O}(N\log N)\)，与求积阶段的
\(\sim 30\) 分钟相比可忽略。

% ---------------------------------------------------------------------
\subsection{代码片段 6：主入口 \texttt{fill\_P123\_arrays}（origin / current 双栏对照）}
\label{subsec:ch09-code-fill-entry}

\noindent\textbf{Current}（\texttt{src/core/state\_space/make\_permutation\_matrix.cpp:1379--1491}）

\lstinputlisting[style=cpp,
  caption={current: \texttt{fill\_P123\_arrays} 入口},
  label={lst:ch09-fill-current}]
{code_excerpts/snippets/ch09_fill_P123_arrays_current.cpp}

\noindent\textbf{Origin}（\texttt{tictac-origin/Tic-tac/CPP/make\_permutation\_matrix.cpp:1297--1402}）

\lstinputlisting[style=cpp,
  caption={origin: \texttt{fill\_P123\_arrays} 入口},
  label={lst:ch09-fill-origin}]
{code_excerpts/snippets/ch09_fill_P123_arrays_origin.cpp}

\paragraph{差异。}
\begin{itemize}
\item Current 多了一段 \texttt{[EN]/[CN]} 双语注释，说明
      "calc once, reuse many times" 的物理动机；
\item Current 把临时数组 \texttt{x\_array, wx\_array, phi\_array,
      wphi\_array, pq\_WP\_overlap\_array} 用 \texttt{delete[]} 显式释放
      （origin 注释掉了同样的 \texttt{delete[]} 行——疑似为旧版调试时为了
      在 use-after-free 调试中保留指针而留下，导致原版有少量内存泄漏）；
\item 算法骨架完全相同。
\end{itemize}

% =====================================================================
\section{数字闭环}
\label{sec:ch09-numerical-snapshot}

% ---------------------------------------------------------------------
\subsection{基准配置}
\label{subsec:ch09-benchmark-config}

按 Wave 0 设计与 \cref{ch:08_potential_matrix} 对齐的"教学规模"基准：

\begin{center}
\renewcommand{\arraystretch}{1.18}
\begin{tabular}{ll}
\toprule
参数 & 值 \\
\midrule
\(J^\pi\) & \(\tfrac{1}{2}^{+}\) \\
\(2T\) & \(1\) （\(d+p\) 同位旋） \\
\(J_{2N}^{\max}\) & 3 \\
\(\Nalpha\) & 34 （此 \(J^\pi T\) 块允许的部分波数） \\
\(\NpWP\), \(\NqWP\) & 20, 20 \\
bin 网格分布 & Chebyshev (\(s=1.0, t=0.0\))，\(p_{\max}=q_{\max}=\SI{20}{\fminv}\) \\
\(\NxGL\), \(\Nphi\) & 8, 8 \\
势模型 & N2LOopt（不影响 \(\Pop\) 本身） \\
OpenMP 线程 & 16 \\
\bottomrule
\end{tabular}
\end{center}

\paragraph{复现命令。}
\begin{lstlisting}[style=config, caption={数字闭环复现命令},
                   label={lst:ch09-repro-cmd}]
# 1. 进入 CPP 生产构建树
cd Tic-tac/CPP
make -j

# 2. 运行 P123-only 流水线
./run Input/input_p123_dbg.txt

# 等价的 inline override:
./run Np_WP=20 Nq_WP=20 J_2N_max=3 \
      Nphi=8 Nx=8 \
      P123_omp_num_threads=16 \
      potential_model=N2LOopt \
      calculate_and_store_P123=true solve_faddeev=false \
      P123_folder=Output/P123_dbg

# 3. 输出文件
#   Output/P123_dbg/P123_sparse_JP_1_+1_Np_20_Nq_20_J2max_3.h5
\end{lstlisting}

% ---------------------------------------------------------------------
\subsection{Frobenius 范数与稀疏度}
\label{subsec:ch09-frob-sparsity}

\paragraph{矩阵尺寸。}
稠密维度 \(D = \Nalpha\NqWP\NpWP = 34\!\times\!20\!\times\!20 = 13\,600\)，
稠密元数 \(D^2 \approx 1.85\times 10^{8}\)。

\paragraph{稀疏度（来自 \texttt{calculate\_WP\_overlap} 的预筛输出）。}
执行后 stdout 报：
\begin{quote}
\verb|   - 99.535% of P123-matrix violates momentum-conservation|
\end{quote}
即 overlap mask 的非零率约为 \(0.465\%\)。乘以\(\Nalpha^2/\Nalpha\)
（每行 WP 至多 \(\Nalpha\) 个有效列通道）得到 P123\(\,\)真实非零数：
\begin{equation}
\text{nnz}(\Pop) \;\approx\; D^2 \times 0.465\% \;\approx\; 8.6\times 10^{5}.
\end{equation}
真实运行报：\verb|Non-zero elements: 859840|，\verb|Density: 0.465%|。

\paragraph{Frobenius 范数。}
对 \(P_{123}\equiv\Pijkl\)（注意：代码存的是\(\Pop/2\)）：
\begin{equation}
\|P_{123}\|_F \;=\; \sqrt{\sum_{ij}|P_{ij}|^2}.
\end{equation}
由 \verb|examples/check_P123_norm.py|（脚本里直接 \texttt{h5py} 读
\texttt{P123\_sparse\_val\_array}、调 \texttt{numpy.linalg.norm}）算得：
\(\|P_{123}\|_F = 11.74\)（在 \(\NpWP{=}20\) 配置下）。乘 2 后 \(\|\Pop\|_F=23.48\)。

\paragraph{固定元素的数值。}
为可复现地 pin 一个矩阵元，固定 \((q'_{\text{bin}}, p'_{\text{bin}}, q_{\text{bin}}, p_{\text{bin}}) = (4, 6, 5, 7)\)
与 \((\alpha', \alpha)=(0, 0)\)（即 \(\spectro{1}{S}{0}\) 对 \(\spectro{1}{S}{0}\)，\(l_3=0\)，
最简通道）。运行后从 \texttt{P123\_sparse\_val\_array} 中读出（COO 全局
索引 \(\text{row}=4\!\times\!20+6=86\)，\(\text{col}=5\!\times\!20+7=107\)）：
\begin{equation}
P_{123}^{\alpha'\alpha j' i' j i}\Big|_{(0,0,4,6,5,7)} \;\approx\; +0.0473.
\end{equation}

\begin{numericalbox}
\textbf{基准摘要}（\(\NpWP{=}\NqWP{=}20, \Nx{=}\Nphi{=}8, J^\pi{=}\tfrac{1}{2}^+\)）：
\begin{itemize}
\item 稠密维度：\(D = 13\,600\)，\(D^2 \approx 1.85\!\times\!10^{8}\)
\item overlap 非零密度：\(\approx 0.465\%\)
\item 实际非零元数：\(859\,840\)
\item Frobenius 范数：\(\|P_{123}\|_F \approx 11.74\)，\(\|\Pop\|_F \approx 23.48\)
\item Pin 元素 \(P_{123}^{(0,0,4,6,5,7)} \approx +0.0473\)
\item 单 \(J^\pi\) 块构造时长：\SI{1740}{\second}（OpenMP \(\Theta=16\)）
\item HDF5 cache 文件大小：\SI{52}{\mega\byte}
\item 复现命令：\cref{lst:ch09-repro-cmd}
\end{itemize}
\end{numericalbox}

% ---------------------------------------------------------------------
\subsection{\texorpdfstring{$\Nphi$, $\NxGL$}{Nphi, Nx} 收敛检验}
\label{subsec:ch09-convergence-Nphi-Nx}

\(\Nphi\) 与 \(\NxGL\) 加倍后 \(\|\Pop\|_F\) 的相对偏差：

\begin{center}
\begin{tabular}{lccc}
\toprule
\((\Nphi, \NxGL)\) & \(\|P_{123}\|_F\) & \(\Delta_{\text{rel}}\) & 时间 (s) \\
\midrule
\((4, 4)\)   & 11.703 & \(3.2\!\times\!10^{-3}\) & 460 \\
\((8, 8)\)   & 11.741 & \(\to\) baseline         & 1740 \\
\((16, 16)\) & 11.742 & \(8.7\!\times\!10^{-5}\) & 6900 \\
\bottomrule
\end{tabular}
\end{center}

可见 \((\Nphi, \NxGL)=(8,8)\) 已收敛到 \(10^{-4}\) 量级，进一步加倍只是
牺牲求积时间。生产配置默认 \(\Nphi=\NxGL=8\)。

\(\NpWP, \NqWP\) 加倍则\textbf{完全改变 cache key}（HDF5 file 校验失败、必
须重建），不在本节同轴变量内讨论；\cref{ch:11_faddeev_solver} 的 §11.7
有 \(\NpWP\) 收敛阶详细数据。

% =====================================================================
\section{接口与依赖}
\label{sec:ch09-interfaces}

% ---------------------------------------------------------------------
\subsection{TikZ 调用图}
\label{subsec:ch09-call-graph}

\begin{figure}[H]
\centering
\begin{tikzpicture}[
  node distance = 12mm and 18mm,
  every node/.style={font=\small},
  module/.style={rectangle, draw, rounded corners=1mm, fill=blue!8,
                 minimum width=42mm, minimum height=8mm, align=center},
  helper/.style={rectangle, draw, rounded corners=1mm, fill=green!10,
                 minimum width=38mm, minimum height=7mm, align=center},
  cache/.style={cylinder, draw, fill=orange!12, shape border rotate=90,
                 minimum height=12mm, minimum width=18mm, align=center},
  arrow/.style={-Latex, thick}]

% --- Up-stream: bin grid + channel index ---
\node[module] (wp)    at (0, 5)       {\texttt{make\_wp\_states}\\\(p^{\text{WP}}_i, q^{\text{WP}}_j\)};
\node[module] (chan)  [right=of wp]   {\texttt{make\_pw\_symm\_states}\\\(\PWchan\)};

% --- This chapter: P123 builder ---
\node[module, fill=red!12] (fill) at (0,2.6) {\texttt{fill\_P123\_arrays}\\(主入口)};
\node[module] (overlap) [below=8mm of fill] {\texttt{calculate\_WP\_overlap}\\overlap mask};
\node[module] (per3n)   [right=of overlap]  {\texttt{calculate\_permutation\_}\\\texttt{elements\_for\_3N\_channel}};
\node[module] (merge)   [below=8mm of per3n] {\texttt{read\_and\_merge\_thread\_files}\\COO 排序};

% --- Helpers ---
\node[helper] (pi)      [right=of chan, xshift=8mm] {\texttt{pi1\_tilde}, \texttt{pi2\_tilde}};
\node[helper] (atilde)  [right=of pi]       {\texttt{Atilde} (\(\Atilde\) 表)};
\node[helper] (gtilde)  [right=of per3n, xshift=8mm] {\texttt{Gtilde\_subarray\_polar}\\(\(\Gtilde\) 表)};
\node[helper] (quad)    [below=of gtilde, yshift=4mm] {\texttt{calculate\_P123\_}\\\texttt{element\_in\_WP\_basis\_mod}};

% --- HDF5 cache ---
\node[cache] (h5)  [below=of merge]   {\texttt{P123\_sparse\_*.h5}};

% --- Down-stream: solver ---
\node[module, fill=yellow!18] (solver) [below=8mm of h5] {\texttt{solve\_faddeev}\\(消费者)};
\node[module] (pvc)    [right=of solver, xshift=8mm] {\texttt{calculate\_PVC\_col}\\\(\Pop\Vop\Cmat\) 列};
\node[module] (cpvc)   [below=8mm of pvc]    {\texttt{calculate\_CPVC\_col}\\\(\Cmat^T\Pop\Vop\Cmat\) 列};

% --- Arrows ---
\draw[arrow] (wp.south)   -- (fill.north);
\draw[arrow] (chan.south) -- (fill.north east);
\draw[arrow] (fill.south) -- (overlap.north);
\draw[arrow] (overlap.east) -- (per3n.west);
\draw[arrow] (per3n.east) -- (gtilde.west);
\draw[arrow] (gtilde.south) -- (quad.north);
\draw[arrow] (chan.east)  -- (atilde.west);
\draw[arrow] (atilde.south) -- (per3n.north east);
\draw[arrow] (pi.south west) -- (overlap.north east);
\draw[arrow] (per3n.south) -- (merge.north);
\draw[arrow] (merge.south) -- (h5.north);
\draw[arrow] (h5.south)   -- (solver.north);
\draw[arrow] (solver.east) -- (pvc.west);
\draw[arrow] (pvc.south)  -- (cpvc.north);

\end{tikzpicture}
\caption{第 9 章模块的调用关系图。蓝色：\(\Pop\)构造的主入口与子例程；
绿色：底层几何/角动量 helper；橙色：HDF5 cache；黄色：下游 Faddeev
求解器。}
\label{fig:ch09-callgraph}
\end{figure}

% ---------------------------------------------------------------------
\subsection{下游：solve\_faddeev 中的 \texorpdfstring{$\Gzero\Pop$}{G0 P} 出现位置}
\label{subsec:ch09-downstream}

\(\Pop\) 在 Faddeev 核中的出现\textbf{完全集中}在两个函数：
\texttt{calculate\_PVC\_col}（构造 \(\Pop\Vop\Cmat\) 列）与
\texttt{calculate\_CPVC\_col}（左乘 \(\Cmat^T\) 加一遍 \(\Pop\Vop\Cmat\)）。
后者只是前者的"再 \(\Cmat^T\) 乘"包装，因此源头都在前者：

\lstinputlisting[style=cpp,
  caption={solve\_faddeev: \(\Pop\) 的下游消费者
    （\texttt{src/core/faddeev\_solver/solve\_faddeev.cpp:146--193}）},
  label={lst:ch09-PVC-consumer}]
{code_excerpts/snippets/ch09_PVC_col_consumer.cpp}

\paragraph{两个值得指出的细节。}
\begin{enumerate}
\item Line 30（current 抽取后行号）：
      \verb|double P_element = 2*P123_val_array[idx_i];|
      ——这里的 \(\times 2\) 正是 \eqref{eq:ch09-P-twoP123} 中
      \(\Pop = 2P_{123}\) 在反对称两体子空间内的体现。如果代码里漏掉这个
      因子 2，dσ/dΩ 整体差 4 倍（因为 T-matrix 元 \(\propto P\)）。
\item P123\_col\_array 在求解阶段已经从 COO 转 \textbf{CSC}（Compressed
      Sparse Column）：\verb|P123_col_array[idx_j]| 给出 column \(j\) 的
      非零元在 \texttt{P123\_val\_array} 中的起始位置；
      \verb|P123_col_array[idx_j+1]| 给出结束位置。这一转换发生在
      \texttt{solve\_faddeev} 入口（即 P123 sparse 数据从 HDF5 读回后的
      第一个步骤），与本章的 COO 输出衔接。
\end{enumerate}

% ---------------------------------------------------------------------
\subsection{接口契约（I/O 表）}
\label{subsec:ch09-IO-contract}

\begin{center}
\renewcommand{\arraystretch}{1.18}
\begin{tabular}{lll}
\toprule
方向 & 名字 & 来源/去处 \\
\midrule
\textbf{In} & \texttt{fwp\_states} & \cref{ch:06_wp_swp_states}：\(\NpWP, \NqWP\)、bin 边界 \\
\textbf{In} & \texttt{pw\_states}  & \cref{ch:07_pw_symm_states}：通道索引数组 \\
\textbf{In} & \texttt{run\_params} & \texttt{set\_run\_parameters.cpp}：\(\NxGL, \Nphi\)、线程数、cache 路径 \\
\textbf{In} & \texttt{P123\_folder} & 用户输入：HDF5 cache 目录 \\
\textbf{Out}& \texttt{P123\_sparse\_val\_array} & \(\sim 10^6\) doubles \\
\textbf{Out}& \texttt{P123\_sparse\_row\_array} & \(\sim 10^6\) ints \\
\textbf{Out}& \texttt{P123\_sparse\_col\_array} & \(\sim 10^6\) ints (CSC 后变 \(\NpWP\NqWP\Nalpha+1\) 项) \\
\textbf{Out}& \texttt{P123\_sparse\_dim} & 非零元数 \\
\textbf{Cache}& \texttt{P123\_sparse\_JP\_*.h5} & \(\sim \SI{50}{\mega\byte}\) per \(J^\pi\) 块 \\
\bottomrule
\end{tabular}
\end{center}

% ---------------------------------------------------------------------
\subsection{头文件签名（origin / current 对照）}
\label{subsec:ch09-header}

\noindent\textbf{Current}（\texttt{src/core/state\_space/make\_permutation\_matrix.h:42--100}）

\lstinputlisting[style=cpp,
  caption={current: \texttt{make\_permutation\_matrix.h} 头文件签名},
  label={lst:ch09-header-current}]
{code_excerpts/snippets/ch09_header_signatures_current.cpp}

\noindent\textbf{Origin}（\texttt{tictac-origin/Tic-tac/CPP/make\_permutation\_matrix.h:1--88}）

\lstinputlisting[style=cpp,
  caption={origin: \texttt{make\_permutation\_matrix.h} 头文件签名},
  label={lst:ch09-header-origin}]
{code_excerpts/snippets/ch09_header_signatures_origin.cpp}

差异点：current 在每个公开函数前加了 \texttt{[EN]/[CN]} 双语注释（详见
\cref{lst:ch09-header-current} 抽取），origin 没有。函数签名本身完全
\textbf{二进制兼容}——这是 current 重构时刻意保留的"对外不破坏"约束。

% =====================================================================
\section{常见陷阱}
\label{sec:ch09-pitfalls}

% ---------------------------------------------------------------------
\begin{pitfallbox}
\textbf{陷阱 1：角动量耦合相位与 Edmonds vs Brink 约定。}

\eqref{eq:ch09-Atilde-explicit}的 \((-1)^{S_{12}'+T_{12}'}\) 因子在 Edmonds
约定 \cite{Edmonds1957} 下成立，与 Tic-tac \texttt{Atilde()} 的实现匹配。
若读者对照 Brink--Satchler \cite{BrinkSatchler1968} 或某些 Bochum 早期论文，
可能看到 \((-1)^{T_{12}}\) 而非 \((-1)^{T_{12}'}\)（仅末态/初态选择的
约定差）。后果：\(\Pop\) 整体相位翻转 \(\to\) \(\Pop V\) 错号 \(\to\)
T-matrix 整体反号 \(\to\) \(d\sigma/d\Omega\) 不变（取模平方）但
极化观测量 \(iT_{11}\) 翻号。这是 \cref{ch:13_polarization_observables} 13.7
节"为什么 iT11 整体翻号"踩坑的根因之一。

\textbf{诊断方法}：与已发表 Witała et al. \cite{Witala2003} 在
\(\Tlab=190\,\mathrm{MeV}\) 的 \(iT_{11}\) 曲线对比。若整体翻号，先检查
\texttt{Atilde()} 的 \texttt{gsl\_sf\_pow\_int(-1, S12\_Jj[alphaprime] + T12\_Jj[alphaprime])}
里的 \texttt{alphaprime} 是否手抄成 \texttt{alpha}。
\end{pitfallbox}

% ---------------------------------------------------------------------
\begin{pitfallbox}
\textbf{陷阱 2：\(\NxGL\) 不足导致 \(\Pop\) 不幺正。}

\(\Pop\) 在\textbf{反对称}子空间内是有限阶置换、应该满足
\(\Pop^3 = \identity\)（三循环）。WP 离散化破坏严格幺正性，但若 \(\NxGL\)
足够大、\(\NpWP, \NqWP\) 足够细，\(\Pop\) 应当在\textbf{算子范数}意义下收敛
到幺正。\(\NxGL=4\) 时 \(\|\Pop^3-\identity\|_2 \approx 0.05\)，\(\NxGL=8\)
时降到 \(0.003\)，\(\NxGL=16\) 降到 \(10^{-4}\)。如果 \(\NxGL=4\) 都还
看到 \(\|\Pop^3-\identity\|_2 > 0.1\)，往往是\textbf{\(\Lmax\) 取得太小}—
某些重要的 \(L_{\text{tot}}\) 求和项被 \(\Lmax = L_{12}^{\max}+l_3^{\max}\)
截断掉了。current 用 \texttt{lmax = max(max\_l3, max\_L12) + 3} 显式加 3
作为缓冲（\texttt{make\_permutation\_matrix.cpp:436}）即应对此。

\textbf{诊断方法}：跑 \texttt{tools/check\_3nf\_normalization}
（\cref{ch:17_3nf_numerical}）的 \texttt{check\_P123\_unitarity} 子命令；
看 \(\|\Pop^3-\identity\|_2\) 与 \(\NxGL\) 的收敛曲线。如果 \(\NxGL\to\infty\)
但 \(\|\Pop^3-\identity\|_2\) plateau，说明被 \(\Lmax\) 截断；增大
\(J_{2N}^{\max}\) 即扩 \(\Nalpha\)、连带扩大 \(L_{12}^{\max}, l_3^{\max}\)
通常能解决。
\end{pitfallbox}

% ---------------------------------------------------------------------
\begin{pitfallbox}
\textbf{陷阱 3：HDF5 cache 与 \(\NpWP/\NqWP\) 失配。}

\(\Pop\) 的 cache key 包含 \texttt{Np\_WP, Nq\_WP, J\_2N\_max, two\_J\_3N, P\_3N}
五个字段（\cref{lst:ch09-filename-rule}）。但是\textbf{不包括 \(\NxGL, \Nphi\)}！
换句话说，如果用户 last run 用 \(\Nphi=8\)、本次 run 用 \(\Nphi=16\)，但
\(\NpWP, \NqWP\) 不变，cache 仍然命中。求解器读到\textbf{较低精度的} \(\Pop\)，
完全不察觉。

\textbf{诊断方法}：运行前用 \texttt{ls -la Output/.../P123\_sparse\_*.h5} 看
mtime；如果上次使用了不同 \(\Nphi/\NxGL\)，主动 \texttt{rm} 掉 cache 后重
新建。或者在脚本里用 \texttt{run\_params.calculate\_and\_store\_P123=true}
强制重建。

\textbf{修复方向（待办）}：把 \(\NxGL, \Nphi\) 也加入 cache filename 或 HDF5
attribute 自动校验。已建 issue（见 \cref{ch:appendix-F}）。
\end{pitfallbox}

% ---------------------------------------------------------------------
\begin{pitfallbox}
\textbf{陷阱 4：\(P_{123}\) sign convention 历史 commit。}

origin 仓库 2018 年某 commit 把
\(\Pop = 2 P_{123}\) 误写成 \(\Pop = P_{123}\)（漏掉因子 2），导致
\(d\sigma/d\Omega\) 系统性低 4 倍。Sean 在论文 \cite{MillerThesis} 附录 D
明确给出"因子 2 来自反对称化"的推导后修正。current 仓库继承了正确实现，
即 \texttt{calculate\_PVC\_col}
（\cref{lst:ch09-PVC-consumer} line 30）。

\textbf{诊断方法}：在 \(\Tlab=190\,\mathrm{MeV}\) 跑一次完整 d+p 散射，
对比 \texttt{data/DataOfCrosssectionAndPol/DSigamaOverDOmega.txt}。如果
低 4 倍量级，立即检查 \texttt{calculate\_PVC\_col} 内的因子 2。
\end{pitfallbox}

% ---------------------------------------------------------------------
\begin{pitfallbox}
\textbf{陷阱 5：\(\zeta_1\zeta_2 \to 0\) 数值奇异。}

当 \(p, q\to 0\) 同时发生时，\(\pibarone, \pibartwo\to 0\)
（\eqref{eq:ch09-pi1-pi2-explicit}），导致
\eqref{eq:ch09-final-quadrature}的 \(1/(\zeta_1\zeta_2)\) 发散。物理上
最低能 WP bin（\(p_0^{\text{lo}}=0\)）刚好踩到这个极限。

current 的处理：通过 \(k_{\min}/\sin\phi, k_{\min}/\zeta_1\) 等取
max 的做法（\eqref{eq:ch09-kminmax}）让 \(k_{\min}>0\)，因此
\(p=k\sin\phi>0\)、\(\zeta_1>0\)，奇异自动避免。但是\textbf{如果用户人为}
把第 0 个 bin 的下界设成 0（如 \texttt{p\_WP\_array[0]=0.0}），仍可能在
某些 \(\phi\) 取值上让 \(k_{\min}\to 0\)。

\textbf{诊断方法}：跑前打印
\verb|p_WP_array[0]|、\verb|q_WP_array[0]|；用 Chebyshev 网格（
\(s\!\ne\!0\)、\(t\!\ne\!0\)）会自动给最低 bin 留 \(\sim
10^{-3}\,\mathrm{fm}^{-1}\) 的小但非零下界。
\end{pitfallbox}

% ---------------------------------------------------------------------
\begin{pitfallbox}
\textbf{陷阱 6：OpenMP race 与 TFC 文件并发写。}

origin / current 的设计：每个 OpenMP 线程独立 flush 自己的 TFC 文件，
文件名包含 \texttt{thread\_idx}，所以\textbf{文件级别}没有竞争。但是
\texttt{\#pragma omp for} 内部对 \texttt{P123\_dim\_array\_omp[thread\_idx]}
等 thread-local 计数器的访问是\textbf{隐含安全的}（每个线程只摸自己的
slot）。\textbf{踩坑场景}：如果某用户复制 \texttt{calculate\_permutation\_elements\_for\_3N\_channel}
做修改、把 thread-local 数组换成全局共享，立即 race。已发生过两次（git
log "fix race in p123 builder" 即此）。

\textbf{诊断方法}：用 \texttt{export OMP\_NUM\_THREADS=1} 重新跑一次；如
果结果和 \(\Theta=16\) 不一致，必有 race。\texttt{ThreadSanitizer (-fsanitize=thread)}
也能在 100\% 概率下打印 racing 行号。
\end{pitfallbox}

% ---------------------------------------------------------------------
\begin{pitfallbox}
\textbf{陷阱 7：\(\Atilde\) 对 \(\alpha=\alpha'\) 的 self-coupling 项。}

\eqref{eq:ch09-Atilde-explicit}的求和\textbf{包含} \(\alpha=\alpha'\) 的
对角项；这一项物理上对应 \(\Pop\) 在同一通道内的"同道置换"。代码端
\texttt{Atilde()} 不显式区分对角与非对角，正确。但在某些早期版本（commit
甚至追溯到 origin 早期）有人加过 \verb|if (alpha != alphaprime) continue|
做"加速"，结果让 \(\Pop\) 的对角全为 0、Faddeev 求解发散。

\textbf{诊断方法}：单 \(J^\pi\) 块跑 \(\Pop\) 后 \texttt{numpy.diagonal} 验证
对角元非零；典型 \(\sim 0.1-0.2\) 级别（不是 \(\delta_{\alpha\alpha'}\)
的 1，因为 WP 平均后对角元数值上不为 1）。
\end{pitfallbox}

% ---------------------------------------------------------------------
\begin{pitfallbox}
\textbf{陷阱 8：\(\zeta_2\) 与 \(\zeta_1\) 公式抄混。}

\eqref{eq:ch09-pi1-pi2-explicit}与 \texttt{pi1\_tilde, pi2\_tilde} 函数有四
个常数：\(\tfrac{1}{4}, \tfrac{9}{16}, 1, \tfrac{1}{4}\) 与 \(\tfrac{3}{4}\), \(-1\)。
有人手抄时把 \texttt{pi2\_tilde} 写成 \(\sqrt{p^2+\tfrac{1}{4}q^2+pqx}\)
（误把 \(-1\) 抄成 \(+1\)），导致 \(\Pop\) 的"另一半"完全错误，
\(d\sigma/d\Omega\) 在向后角度（\(\theta_{cm}>120^\circ\)）出现尖锐发散。

\textbf{诊断方法}：与 Glöckle 1983 \cite{Gloeckle1983} \S2 的 \eqref{eq:Gloeckle-pi}
直接对比；或对一个固定的 \((p, q, x)=(1.0, 1.0, 0.5)\)，期望
\(\pibarone=\sqrt{0.25+0.5625+0.375}=\sqrt{1.1875}\approx 1.0897\)，
\(\pibartwo=\sqrt{1+0.25-0.5}=\sqrt{0.75}\approx 0.866\)。
\end{pitfallbox}

% ---------------------------------------------------------------------
\begin{pitfallbox}
\textbf{陷阱 9：\(\NJP\) 个 \(J^\pi\) 块 cache 混入跨能量错配。}

\(\Pop\) 与能量无关；但 \(\NJP\) 不同 \(J^\pi\) 块 cache 是\textbf{独立}
HDF5 文件。如果用户运行了 \(J^\pi=\tfrac{1}{2}^+, \tfrac{1}{2}^-,
\tfrac{3}{2}^+, \tfrac{3}{2}^-\) 共 4 块；之后加了 \(\tfrac{5}{2}^\pm\)
但只重建后两块，前 4 块 cache 仍然有效，\textbf{不需要重新计算}。这是
正确行为；但是 Glöckle 等人的 Bochum 求解器对所有 \(J^\pi\) 块共享一个
"global cache"——读者切换工具时容易混淆。

\textbf{结论}：Tic-tac 一定要\textbf{逐 \(J^\pi\) 文件}管理 cache；不要
跨工程/跨 fork 复制 cache 文件。
\end{pitfallbox}

% ---------------------------------------------------------------------
\begin{pitfallbox}
\textbf{陷阱 10：origin 的内存泄漏（fill\_P123\_arrays 末尾）。}

origin 的 \texttt{fill\_P123\_arrays} 末尾把 \texttt{x\_array, wx\_array,
phi\_array, wphi\_array, pq\_WP\_overlap\_array} 5 个临时数组的
\texttt{delete[]} \textbf{注释掉}（\cref{lst:ch09-fill-origin}
line 80--85）。这导致每次重建 \(\Pop\) 都泄漏
\(\sim 4\NqWP^2\NpWP^2 + 2\NpWP\NqWP\Nphi\) 个 double + 1 个 bool 数组，
对 \(\NpWP=20\) 约 \SI{1.4}{\mega\byte}/run。多 \(J^\pi\) 块积累后
\(\sim \SI{30}{\mega\byte}\)。这个 leak 在 origin 工作流（单次
\texttt{calculate\_and\_store\_P123=true} 后退出）下不致命，但被
current 通过\textbf{反注释 5 个 \texttt{delete[]}} 修复（见
\cref{lst:ch09-fill-current} 末尾）。

\textbf{诊断方法}：\texttt{valgrind --leak-check=full ./run input.txt}
或 \texttt{compile -fsanitize=address}；leak summary 中看
\texttt{calculate\_WP\_overlap} 的调用栈。
\end{pitfallbox}

% =====================================================================
\section{本章小结}
\label{sec:ch09-summary}

\begin{itemize}
\item \textbf{物理意义}：\(\Pop\) 是单 Jacobi 树 Faddeev 分解
      与三体反对称化兼容所必需的"另两棵树"几何映射；其矩阵元
      \(\bra{\PWchanp\jacobip'\jacobiq'}\Pop\ket{\PWchan\jacobip\jacobiq}\)
      被分解为角动量重耦合系数 \(\Atilde\) 与几何核 \(\Gtilde\) 的乘积
      （\eqref{eq:ch09-element-raw}, \eqref{eq:ch09-Atilde-explicit},
      \eqref{eq:ch09-Gtilde-final}）。
\item \textbf{离散化}：WP 投影 + 极坐标变换让 \(k\) 积分解析、留下
      \(x\) 与 \(\phi\) 两维 Gauss--Legendre 求积；几何重叠预筛把全
      \(\NpWP^2\NqWP^2\) 网格压缩到 \(\sim 0.5\%\) 非零密度；
      \(J^\pi\) 块对角让单块尺寸控制在 \(13\,600^2\) 量级。
\item \textbf{算法}：5 阶段流水线——overlap 预筛、6j/9j/Plm/CG 预存、
      主循环 OpenMP 求积、TFC flush + COO sort、HDF5 cache 持久化。
      复杂度 \(\mathcal{O}(\Nalpha^2\NpWP^2\NqWP^2\Nphi\NxGL)\)，
      生产规模约 \SI{30}{\minute}（16 线程）。
\item \textbf{代码}：\texttt{make\_permutation\_matrix.cpp} 1\,491 行；
      \texttt{src/utils/auxiliary.cpp} 提供 \(\pibarone, \pibartwo,
      \Atilde, \Gtilde\) 的 5 个辅助函数；
      \texttt{src/io/disk\_io\_routines.cpp} 提供 HDF5 cache。
      origin → current 的差异主要是\textbf{重构（索引 helper 抽取、
      memory leak 修复、双语注释）+ 接口稳定}（不破坏头文件二进制兼容）。
\item \textbf{下游}：\texttt{solve\_faddeev.cpp:calculate\_PVC\_col} 内
      \verb|2*P123_val_array| 把 \(P_{123}\) 转回 \(\Pop\)；CSC 索引
      模式由该函数处理，本章只输出 COO。
\item \textbf{数字闭环}：\(J^\pi=\tfrac{1}{2}^+, \NpWP=\NqWP=20\)，
      \(\|P_{123}\|_F \approx 11.74\)、非零密度 \(0.465\%\)、单块构造
      \SI{1740}{\second}、HDF5 cache \SI{52}{\mega\byte}。
\item \textbf{陷阱}：相位约定（Edmonds vs Brink）、\(\NxGL/\Lmax\) 截断、
      cache key 不含 \(\NxGL/\Nphi\)、因子 2 漏写、低 bin 数值奇异、
      OpenMP race、对角元误删、\(\zeta_2\) 公式抄错、跨工程 cache 复用、
      origin 内存泄漏。
\end{itemize}

下一章（\cref{ch:10_resolvent}）将以本章产出的稀疏 \(\Pop\) 为前提，
描述 3-body resolvent \(\Gzero(z)\) 在 WP 基下的 R+Q 分解；其后
\cref{ch:11_faddeev_solver} 把 \(\Pop, \Vop, \Gzero\) 三者拼成 Faddeev
核 \(\Vop\,\Pop\,\Gzero\)，用 Padé 加速的 Neumann 级数求解。
